% Electromagnetic Wave Propagation Template
% Topics: Maxwell's equations, Fresnel coefficients, waveguides, skin depth
% Style: Technical report with computational electromagnetics analysis

\documentclass[a4paper, 11pt]{article}
\usepackage[utf8]{inputenc}
\usepackage[T1]{fontenc}
\usepackage{amsmath, amssymb}
\usepackage{graphicx}
\usepackage{siunitx}
\usepackage{booktabs}
\usepackage{subcaption}
\usepackage[makestderr]{pythontex}

% Theorem environments
\newtheorem{definition}{Definition}[section]
\newtheorem{theorem}{Theorem}[section]
\newtheorem{example}{Example}[section]
\newtheorem{remark}{Remark}[section]

\title{Electromagnetic Wave Propagation: From Maxwell's Equations to Waveguides}
\author{Electromagnetics Research Laboratory}
\date{\today}

\begin{document}
\maketitle

\begin{abstract}
This technical report presents a comprehensive computational analysis of electromagnetic
wave propagation phenomena derived from Maxwell's equations. We examine plane wave
solutions in free space and lossy media, analyze reflection and transmission at dielectric
interfaces using Fresnel coefficients, investigate rectangular waveguide modes with cutoff
frequencies and dispersion relations, and characterize skin depth and attenuation in
conductors. The analysis includes numerical computation of Brewster angles, total internal
reflection conditions, TE/TM mode field distributions, and frequency-dependent penetration
depths relevant to RF engineering and microwave systems.
\end{abstract}

\section{Introduction}

Electromagnetic wave propagation forms the foundation of modern wireless communication,
radar systems, optical fibers, and microwave engineering. Maxwell's equations, formulated
in 1865, unify electricity and magnetism and predict the existence of electromagnetic waves
traveling at the speed of light.

\begin{definition}[Maxwell's Equations in Source-Free Media]
In regions without charges or currents, Maxwell's equations are:
\begin{align}
\nabla \times \mathbf{E} &= -\frac{\partial \mathbf{B}}{\partial t} \quad \text{(Faraday's law)}\\
\nabla \times \mathbf{H} &= \frac{\partial \mathbf{D}}{\partial t} \quad \text{(Ampère-Maxwell law)}\\
\nabla \cdot \mathbf{D} &= 0 \quad \text{(Gauss's law)}\\
\nabla \cdot \mathbf{B} &= 0 \quad \text{(No magnetic monopoles)}
\end{align}
where $\mathbf{D} = \varepsilon \mathbf{E}$ and $\mathbf{B} = \mu \mathbf{H}$.
\end{definition}

\section{Theoretical Framework}

\subsection{Wave Equation and Plane Waves}

\begin{theorem}[Wave Equation]
Combining Faraday's and Ampère's laws yields the wave equation:
\begin{equation}
\nabla^2 \mathbf{E} = \mu \varepsilon \frac{\partial^2 \mathbf{E}}{\partial t^2}
\end{equation}
with phase velocity $v_p = 1/\sqrt{\mu \varepsilon}$ and in vacuum $c = 1/\sqrt{\mu_0 \varepsilon_0}$.
\end{theorem}

\begin{definition}[Plane Wave Solution]
A uniform plane wave propagating in the $+z$ direction has the form:
\begin{equation}
\mathbf{E}(z,t) = E_0 \cos(\omega t - \beta z + \phi) \hat{\mathbf{x}}
\end{equation}
where $\beta = \omega\sqrt{\mu\varepsilon}$ is the phase constant and the magnetic field is:
\begin{equation}
\mathbf{H}(z,t) = \frac{E_0}{\eta} \cos(\omega t - \beta z + \phi) \hat{\mathbf{y}}
\end{equation}
with intrinsic impedance $\eta = \sqrt{\mu/\varepsilon}$ (377 $\Omega$ in free space).
\end{definition}

\subsection{Reflection and Transmission at Interfaces}

\begin{theorem}[Fresnel Equations]
For a plane wave incident at angle $\theta_i$ on a dielectric interface, the reflection
and transmission coefficients for perpendicular (TE) polarization are:
\begin{align}
\Gamma_\perp &= \frac{\eta_2 \cos\theta_i - \eta_1 \cos\theta_t}{\eta_2 \cos\theta_i + \eta_1 \cos\theta_t}\\
\tau_\perp &= \frac{2\eta_2 \cos\theta_i}{\eta_2 \cos\theta_i + \eta_1 \cos\theta_t}
\end{align}
For parallel (TM) polarization:
\begin{align}
\Gamma_\parallel &= \frac{\eta_2 \cos\theta_t - \eta_1 \cos\theta_i}{\eta_2 \cos\theta_t + \eta_1 \cos\theta_i}\\
\tau_\parallel &= \frac{2\eta_2 \cos\theta_i}{\eta_2 \cos\theta_t + \eta_1 \cos\theta_i}
\end{align}
where $\theta_t$ is determined by Snell's law: $n_1 \sin\theta_i = n_2 \sin\theta_t$.
\end{theorem}

\begin{definition}[Brewster Angle]
The Brewster angle $\theta_B$ is the incidence angle at which parallel polarization
has zero reflection:
\begin{equation}
\theta_B = \arctan\left(\frac{n_2}{n_1}\right)
\end{equation}
At this angle, reflected and transmitted rays are perpendicular.
\end{definition}

\subsection{Propagation in Lossy Media}

\begin{theorem}[Complex Propagation Constant]
In a lossy medium with conductivity $\sigma$, the propagation constant is complex:
\begin{equation}
\gamma = \alpha + j\beta = j\omega\sqrt{\mu\varepsilon_c} = j\omega\sqrt{\mu\varepsilon\left(1 - j\frac{\sigma}{\omega\varepsilon}\right)}
\end{equation}
where $\alpha$ is the attenuation constant (Np/m) and $\beta$ is the phase constant (rad/m).
\end{theorem}

\begin{definition}[Skin Depth]
The skin depth $\delta$ is the distance over which field amplitude decays to $1/e$:
\begin{equation}
\delta = \frac{1}{\alpha} = \sqrt{\frac{2}{\omega\mu\sigma}}
\end{equation}
For good conductors ($\sigma \gg \omega\varepsilon$), this simplifies to the classical formula.
\end{definition}

\subsection{Rectangular Waveguide Theory}

\begin{theorem}[Cutoff Frequency]
For a rectangular waveguide with dimensions $a \times b$ (where $a > b$), the cutoff
frequency for the $\text{TE}_{mn}$ or $\text{TM}_{mn}$ mode is:
\begin{equation}
f_{c,mn} = \frac{c}{2\sqrt{\mu_r\varepsilon_r}} \sqrt{\left(\frac{m}{a}\right)^2 + \left(\frac{n}{b}\right)^2}
\end{equation}
The dominant mode is $\text{TE}_{10}$ with $f_{c,10} = c/(2a\sqrt{\mu_r\varepsilon_r})$.
\end{theorem}

\begin{definition}[Phase and Group Velocity]
Above cutoff, the phase velocity is:
\begin{equation}
v_p = \frac{c}{\sqrt{1 - (f_c/f)^2}}
\end{equation}
and the group velocity (energy propagation speed) is:
\begin{equation}
v_g = c\sqrt{1 - (f_c/f)^2}
\end{equation}
Note that $v_p v_g = c^2$.
\end{definition}

\section{Computational Analysis}

\begin{pycode}
import numpy as np
import matplotlib.pyplot as plt
from scipy.constants import c, mu_0, epsilon_0
from scipy.optimize import fsolve

np.random.seed(42)

# Physical constants
eta_0 = np.sqrt(mu_0 / epsilon_0)  # Free space impedance (377 ohms)

# Define material properties
def refractive_index(epsilon_r, mu_r=1.0):
    return np.sqrt(epsilon_r * mu_r)

def intrinsic_impedance(epsilon_r, mu_r=1.0):
    return eta_0 * np.sqrt(mu_r / epsilon_r)

# Fresnel coefficients
def fresnel_perpendicular(theta_i, n1, n2):
    """TE polarization (s-polarized)"""
    sin_theta_t = (n1 / n2) * np.sin(theta_i)
    if np.any(sin_theta_t > 1):
        # Total internal reflection
        cos_theta_t = 1j * np.sqrt(sin_theta_t**2 - 1)
    else:
        cos_theta_t = np.sqrt(1 - sin_theta_t**2)

    cos_theta_i = np.cos(theta_i)
    Gamma = (n1 * cos_theta_i - n2 * cos_theta_t) / (n1 * cos_theta_i + n2 * cos_theta_t)
    tau = (2 * n1 * cos_theta_i) / (n1 * cos_theta_i + n2 * cos_theta_t)
    return Gamma, tau

def fresnel_parallel(theta_i, n1, n2):
    """TM polarization (p-polarized)"""
    sin_theta_t = (n1 / n2) * np.sin(theta_i)
    if np.any(sin_theta_t > 1):
        cos_theta_t = 1j * np.sqrt(sin_theta_t**2 - 1)
    else:
        cos_theta_t = np.sqrt(1 - sin_theta_t**2)

    cos_theta_i = np.cos(theta_i)
    Gamma = (n2 * cos_theta_i - n1 * cos_theta_t) / (n2 * cos_theta_i + n1 * cos_theta_t)
    tau = (2 * n1 * cos_theta_i) / (n2 * cos_theta_i + n1 * cos_theta_t)
    return Gamma, tau

def brewster_angle(n1, n2):
    return np.arctan(n2 / n1)

def critical_angle(n1, n2):
    if n1 <= n2:
        return None
    return np.arcsin(n2 / n1)

# Skin depth calculation
def skin_depth(frequency, sigma, mu_r=1.0):
    """Skin depth in meters"""
    omega = 2 * np.pi * frequency
    return np.sqrt(2 / (omega * mu_0 * mu_r * sigma))

def attenuation_constant(frequency, sigma, epsilon_r=1.0, mu_r=1.0):
    """Attenuation constant in Np/m"""
    omega = 2 * np.pi * frequency
    eps = epsilon_0 * epsilon_r
    mu = mu_0 * mu_r
    # General formula for lossy dielectric
    loss_tangent = sigma / (omega * eps)
    beta_lossless = omega * np.sqrt(mu * eps)
    alpha = beta_lossless * loss_tangent / 2
    return alpha

# Waveguide properties
def cutoff_frequency(m, n, a, b, epsilon_r=1.0, mu_r=1.0):
    """Cutoff frequency for TE_mn or TM_mn mode"""
    return (c / (2 * np.sqrt(epsilon_r * mu_r))) * np.sqrt((m/a)**2 + (n/b)**2)

def phase_velocity(frequency, cutoff_freq):
    """Phase velocity in waveguide"""
    if frequency <= cutoff_freq:
        return np.inf
    return c / np.sqrt(1 - (cutoff_freq / frequency)**2)

def group_velocity(frequency, cutoff_freq):
    """Group velocity in waveguide"""
    if frequency <= cutoff_freq:
        return 0
    return c * np.sqrt(1 - (cutoff_freq / frequency)**2)

def guide_wavelength(frequency, cutoff_freq):
    """Wavelength in waveguide"""
    lambda_0 = c / frequency
    if frequency <= cutoff_freq:
        return np.inf
    return lambda_0 / np.sqrt(1 - (cutoff_freq / frequency)**2)

# Materials
n_air = 1.0
n_glass = 1.5
n_water = 1.33
epsilon_glass = n_glass**2
epsilon_water = n_water**2

# Compute Brewster angle
theta_b_air_glass = brewster_angle(n_air, n_glass)
theta_b_air_water = brewster_angle(n_air, n_water)

# Compute critical angles
theta_c_glass_air = critical_angle(n_glass, n_air)
theta_c_water_air = critical_angle(n_water, n_air)

# Create comprehensive figure
fig = plt.figure(figsize=(16, 14))

# Plot 1: Fresnel coefficients vs angle (air to glass)
ax1 = fig.add_subplot(3, 3, 1)
angles_deg = np.linspace(0, 89, 500)
angles_rad = np.radians(angles_deg)

Gamma_perp_ag, tau_perp_ag = fresnel_perpendicular(angles_rad, n_air, n_glass)
Gamma_para_ag, tau_para_ag = fresnel_parallel(angles_rad, n_air, n_glass)

ax1.plot(angles_deg, np.abs(Gamma_perp_ag)**2, 'b-', linewidth=2, label='$R_\\perp$ (TE)')
ax1.plot(angles_deg, np.abs(Gamma_para_ag)**2, 'r-', linewidth=2, label='$R_\\parallel$ (TM)')
ax1.axvline(x=np.degrees(theta_b_air_glass), color='gray', linestyle='--',
            alpha=0.7, label=f'Brewster: {np.degrees(theta_b_air_glass):.1f}°')
ax1.set_xlabel('Incidence Angle (degrees)')
ax1.set_ylabel('Reflectance')
ax1.set_title('Fresnel Reflectance: Air $\\to$ Glass')
ax1.legend(fontsize=9)
ax1.grid(True, alpha=0.3)
ax1.set_xlim(0, 90)
ax1.set_ylim(0, 1)

# Plot 2: Transmittance vs angle
ax2 = fig.add_subplot(3, 3, 2)
T_perp = 1 - np.abs(Gamma_perp_ag)**2
T_para = 1 - np.abs(Gamma_para_ag)**2

ax2.plot(angles_deg, T_perp, 'b-', linewidth=2, label='$T_\\perp$ (TE)')
ax2.plot(angles_deg, T_para, 'r-', linewidth=2, label='$T_\\parallel$ (TM)')
ax2.axvline(x=np.degrees(theta_b_air_glass), color='gray', linestyle='--', alpha=0.7)
ax2.set_xlabel('Incidence Angle (degrees)')
ax2.set_ylabel('Transmittance')
ax2.set_title('Fresnel Transmittance: Air $\\to$ Glass')
ax2.legend(fontsize=9)
ax2.grid(True, alpha=0.3)
ax2.set_xlim(0, 90)
ax2.set_ylim(0, 1)

# Plot 3: Total internal reflection (glass to air)
ax3 = fig.add_subplot(3, 3, 3)
angles_deg_ga = np.linspace(0, 89, 500)
angles_rad_ga = np.radians(angles_deg_ga)

Gamma_perp_ga, _ = fresnel_perpendicular(angles_rad_ga, n_glass, n_air)
Gamma_para_ga, _ = fresnel_parallel(angles_rad_ga, n_glass, n_air)

ax3.plot(angles_deg_ga, np.abs(Gamma_perp_ga)**2, 'b-', linewidth=2, label='$R_\\perp$')
ax3.plot(angles_deg_ga, np.abs(Gamma_para_ga)**2, 'r-', linewidth=2, label='$R_\\parallel$')
ax3.axvline(x=np.degrees(theta_c_glass_air), color='green', linestyle='--',
            linewidth=2, label=f'Critical: {np.degrees(theta_c_glass_air):.1f}°')
ax3.set_xlabel('Incidence Angle (degrees)')
ax3.set_ylabel('Reflectance')
ax3.set_title('Total Internal Reflection: Glass $\\to$ Air')
ax3.legend(fontsize=9)
ax3.grid(True, alpha=0.3)
ax3.set_xlim(0, 90)
ax3.set_ylim(0, 1.05)

# Plot 4: Skin depth vs frequency in copper
ax4 = fig.add_subplot(3, 3, 4)
sigma_copper = 5.96e7  # S/m
frequencies_khz = np.logspace(1, 4, 100)  # 10 kHz to 10 MHz
frequencies_hz = frequencies_khz * 1e3

skin_depths_copper = skin_depth(frequencies_hz, sigma_copper)

ax4.loglog(frequencies_khz, skin_depths_copper * 1e6, 'b-', linewidth=2)
ax4.set_xlabel('Frequency (kHz)')
ax4.set_ylabel('Skin Depth ($\\mu$m)')
ax4.set_title('Skin Depth in Copper ($\\sigma$ = 5.96$\\times$10$^7$ S/m)')
ax4.grid(True, alpha=0.3, which='both')
ax4.axhline(y=10, color='gray', linestyle='--', alpha=0.5, label='10 $\\mu$m')
ax4.axhline(y=100, color='gray', linestyle=':', alpha=0.5, label='100 $\\mu$m')
ax4.legend(fontsize=8)

# Plot 5: Skin depth comparison for different metals
ax5 = fig.add_subplot(3, 3, 5)
sigma_aluminum = 3.5e7  # S/m
sigma_gold = 4.1e7  # S/m
sigma_silver = 6.3e7  # S/m

freq_mhz = np.logspace(-1, 2, 100)  # 0.1 MHz to 100 MHz
freq_hz_metal = freq_mhz * 1e6

skin_copper = skin_depth(freq_hz_metal, sigma_copper) * 1e6
skin_aluminum = skin_depth(freq_hz_metal, sigma_aluminum) * 1e6
skin_gold = skin_depth(freq_hz_metal, sigma_gold) * 1e6
skin_silver = skin_depth(freq_hz_metal, sigma_silver) * 1e6

ax5.loglog(freq_mhz, skin_copper, 'b-', linewidth=2, label='Copper')
ax5.loglog(freq_mhz, skin_aluminum, 'r-', linewidth=2, label='Aluminum')
ax5.loglog(freq_mhz, skin_gold, 'g-', linewidth=2, label='Gold')
ax5.loglog(freq_mhz, skin_silver, 'm-', linewidth=2, label='Silver')
ax5.set_xlabel('Frequency (MHz)')
ax5.set_ylabel('Skin Depth ($\\mu$m)')
ax5.set_title('Skin Depth: Metal Comparison')
ax5.legend(fontsize=9)
ax5.grid(True, alpha=0.3, which='both')

# Plot 6: Waveguide dispersion (TE10 mode)
ax6 = fig.add_subplot(3, 3, 6)
a_guide = 0.02286  # Standard WR-90 waveguide (m)
b_guide = 0.01016  # (m)

fc_10 = cutoff_frequency(1, 0, a_guide, b_guide)
fc_20 = cutoff_frequency(2, 0, a_guide, b_guide)
fc_01 = cutoff_frequency(0, 1, a_guide, b_guide)
fc_11 = cutoff_frequency(1, 1, a_guide, b_guide)

freq_ghz = np.linspace(fc_10 / 1e9, 2 * fc_20 / 1e9, 500)
freq_hz_guide = freq_ghz * 1e9

vp_10 = np.array([phase_velocity(f, fc_10) for f in freq_hz_guide])
vg_10 = np.array([group_velocity(f, fc_10) for f in freq_hz_guide])

# Normalize to speed of light
vp_norm = vp_10 / c
vg_norm = vg_10 / c

ax6.plot(freq_ghz, vp_norm, 'b-', linewidth=2, label='$v_p/c$ (phase)')
ax6.plot(freq_ghz, vg_norm, 'r-', linewidth=2, label='$v_g/c$ (group)')
ax6.axvline(x=fc_10/1e9, color='gray', linestyle='--', alpha=0.7,
            label=f'$f_{{c,10}}$ = {fc_10/1e9:.2f} GHz')
ax6.axhline(y=1, color='black', linestyle=':', alpha=0.5)
ax6.set_xlabel('Frequency (GHz)')
ax6.set_ylabel('Normalized Velocity')
ax6.set_title('Waveguide Dispersion: TE$_{10}$ Mode (WR-90)')
ax6.legend(fontsize=9)
ax6.grid(True, alpha=0.3)
ax6.set_xlim(fc_10/1e9, 2*fc_20/1e9)
ax6.set_ylim(0, 3)

# Plot 7: Multiple mode cutoff frequencies
ax7 = fig.add_subplot(3, 3, 7)
modes = ['TE$_{10}$', 'TE$_{20}$', 'TE$_{01}$', 'TE$_{11}$', 'TM$_{11}$']
cutoffs_ghz = [fc_10/1e9, fc_20/1e9, fc_01/1e9, fc_11/1e9, fc_11/1e9]

y_pos = np.arange(len(modes))
bars = ax7.barh(y_pos, cutoffs_ghz, color='steelblue', edgecolor='black')
bars[0].set_color('coral')  # Highlight dominant mode
ax7.set_yticks(y_pos)
ax7.set_yticklabels(modes, fontsize=10)
ax7.set_xlabel('Cutoff Frequency (GHz)')
ax7.set_title('Mode Cutoff Frequencies (WR-90)')
ax7.grid(True, alpha=0.3, axis='x')

# Add frequency labels
for i, (bar, fc) in enumerate(zip(bars, cutoffs_ghz)):
    ax7.text(fc + 0.5, i, f'{fc:.2f} GHz', va='center', fontsize=9)

# Plot 8: TE10 electric field distribution
ax8 = fig.add_subplot(3, 3, 8)
x_norm = np.linspace(0, 1, 100)
y_norm = np.linspace(0, 1, 100)
X_norm, Y_norm = np.meshgrid(x_norm, y_norm)

# TE10 mode: Ey component (dominant)
E_field_te10 = np.sin(np.pi * X_norm)

im = ax8.contourf(X_norm, Y_norm, E_field_te10, levels=20, cmap='RdBu_r')
plt.colorbar(im, ax=ax8, label='$E_y$ (normalized)')
ax8.set_xlabel('x/a')
ax8.set_ylabel('y/b')
ax8.set_title('TE$_{10}$ Mode: Electric Field ($E_y$)')
ax8.set_aspect('equal')

# Plot 9: Guide wavelength vs frequency
ax9 = fig.add_subplot(3, 3, 9)
lambda_guide = np.array([guide_wavelength(f, fc_10) for f in freq_hz_guide])
lambda_free = c / freq_hz_guide

ax9.plot(freq_ghz, lambda_free * 1000, 'b--', linewidth=2, label='Free space ($\\lambda_0$)')
ax9.plot(freq_ghz, lambda_guide * 1000, 'r-', linewidth=2, label='Waveguide ($\\lambda_g$)')
ax9.axvline(x=fc_10/1e9, color='gray', linestyle='--', alpha=0.7)
ax9.set_xlabel('Frequency (GHz)')
ax9.set_ylabel('Wavelength (mm)')
ax9.set_title('Wavelength Comparison')
ax9.legend(fontsize=9)
ax9.grid(True, alpha=0.3)
ax9.set_xlim(fc_10/1e9, 2*fc_20/1e9)
ax9.set_ylim(0, 100)

plt.tight_layout()
plt.savefig('wave_propagation_analysis.pdf', dpi=150, bbox_inches='tight')
plt.close()
\end{pycode}

\begin{figure}[htbp]
\centering
\includegraphics[width=\textwidth]{wave_propagation_analysis.pdf}
\caption{Comprehensive electromagnetic wave propagation analysis: (a) Fresnel reflectance
coefficients for air-to-glass interface showing Brewster angle at 56.3° where parallel
polarization has zero reflection; (b) Corresponding transmittance revealing complementary
behavior; (c) Total internal reflection at glass-to-air interface with critical angle
at 41.8° beyond which complete reflection occurs; (d) Skin depth in copper decreasing
from 660 $\mu$m at 10 kHz to 21 $\mu$m at 1 MHz; (e) Comparative skin depths for
Cu, Al, Au, and Ag showing frequency-dependent penetration; (f) Phase and group velocity
dispersion in WR-90 waveguide for TE$_{10}$ mode with cutoff at 6.56 GHz; (g) Cutoff
frequencies for multiple waveguide modes; (h) TE$_{10}$ mode electric field distribution
exhibiting half-wave sinusoidal variation; (i) Guide wavelength exceeding free-space
wavelength near cutoff frequency.}
\label{fig:propagation}
\end{figure}

\section{Results}

\subsection{Interface Phenomena}

\begin{pycode}
print(r"\begin{table}[htbp]")
print(r"\centering")
print(r"\caption{Critical Angles for Common Material Interfaces}")
print(r"\begin{tabular}{lccc}")
print(r"\toprule")
print(r"Interface & $n_1$ & $n_2$ & Brewster Angle & Critical Angle \\")
print(r"\midrule")

interfaces = [
    ("Air $\\to$ Glass", n_air, n_glass, theta_b_air_glass, None),
    ("Air $\\to$ Water", n_air, n_water, theta_b_air_water, None),
    ("Glass $\\to$ Air", n_glass, n_air, brewster_angle(n_glass, n_air), theta_c_glass_air),
    ("Water $\\to$ Air", n_water, n_air, brewster_angle(n_water, n_air), theta_c_water_air),
]

for name, n1, n2, theta_b, theta_c in interfaces:
    theta_b_deg = np.degrees(theta_b)
    if theta_c is not None:
        theta_c_deg = np.degrees(theta_c)
        print(f"{name} & {n1:.2f} & {n2:.2f} & {theta_b_deg:.1f}$^\\circ$ & {theta_c_deg:.1f}$^\\circ$ \\\\")
    else:
        print(f"{name} & {n1:.2f} & {n2:.2f} & {theta_b_deg:.1f}$^\\circ$ & --- \\\\")

print(r"\bottomrule")
print(r"\end{tabular}")
print(r"\label{tab:angles}")
print(r"\end{table}")
\end{pycode}

\subsection{Skin Depth Analysis}

\begin{pycode}
print(r"\begin{table}[htbp]")
print(r"\centering")
print(r"\caption{Skin Depth in Copper at Representative Frequencies}")
print(r"\begin{tabular}{cccc}")
print(r"\toprule")
print(r"Frequency & Wavelength & Skin Depth & Attenuation \\")
print(r"& (m) & ($\mu$m) & (dB/cm) \\")
print(r"\midrule")

test_frequencies = [60, 1e3, 1e6, 1e9, 10e9]  # Hz

for f in test_frequencies:
    wavelength = c / f
    delta = skin_depth(f, sigma_copper) * 1e6  # Convert to micrometers
    alpha_np = 1 / (delta * 1e-6)  # Np/m
    alpha_db = alpha_np * 8.686  # dB/m

    if f < 1e3:
        f_str = f"{f:.0f} Hz"
    elif f < 1e6:
        f_str = f"{f/1e3:.0f} kHz"
    elif f < 1e9:
        f_str = f"{f/1e6:.0f} MHz"
    else:
        f_str = f"{f/1e9:.0f} GHz"

    if wavelength >= 1:
        wl_str = f"{wavelength:.1f}"
    elif wavelength >= 0.001:
        wl_str = f"{wavelength*1000:.1f}$\\times$10$^{{-3}}$"
    else:
        wl_str = f"{wavelength:.2e}"

    print(f"{f_str} & {wl_str} & {delta:.1f} & {alpha_db/100:.2f} \\\\")

print(r"\bottomrule")
print(r"\end{tabular}")
print(r"\label{tab:skin}")
print(r"\end{table}")
\end{pycode}

\subsection{Waveguide Characteristics}

\begin{pycode}
print(r"\begin{table}[htbp]")
print(r"\centering")
print(r"\caption{Standard Rectangular Waveguide Parameters}")
print(r"\begin{tabular}{lccccc}")
print(r"\toprule")
print(r"Designation & $a$ (mm) & $b$ (mm) & $f_{c,10}$ (GHz) & Band & Frequency Range \\")
print(r"\midrule")

waveguides = [
    ("WR-284", 72.14, 34.04, "WG-6", "2.6--3.95 GHz"),
    ("WR-187", 47.55, 22.15, "WG-9", "3.95--5.85 GHz"),
    ("WR-90", 22.86, 10.16, "X-band", "8.2--12.4 GHz"),
    ("WR-62", 15.80, 7.90, "Ku-band", "12.4--18.0 GHz"),
    ("WR-42", 10.67, 4.32, "Ka-band", "18.0--26.5 GHz"),
]

for name, a_mm, b_mm, band, freq_range in waveguides:
    a_m = a_mm / 1000
    b_m = b_mm / 1000
    fc = cutoff_frequency(1, 0, a_m, b_m) / 1e9
    print(f"{name} & {a_mm:.2f} & {b_mm:.2f} & {fc:.2f} & {band} & {freq_range} \\\\")

print(r"\bottomrule")
print(r"\end{tabular}")
print(r"\label{tab:waveguide}")
print(r"\end{table}")
\end{pycode}

\section{Discussion}

\begin{example}[Polarization-Dependent Reflection]
At the air-glass Brewster angle ($56.3°$), parallel-polarized light experiences zero
reflection while perpendicular-polarized light has approximately 7\% reflectance. This
phenomenon enables polarizing beam splitters and laser cavity windows oriented at
Brewster's angle to minimize losses for one polarization.
\end{example}

\begin{example}[Total Internal Reflection Applications]
The critical angle for glass-to-air ($41.8°$) enables optical fiber communication.
Light launched into a fiber core at angles exceeding the critical angle undergoes
total internal reflection, confining the signal over kilometer-scale distances with
minimal loss. This principle also governs prism reflectors and retroreflectors.
\end{example}

\begin{remark}[RF Shielding and Skin Effect]
At 1 MHz, the skin depth in copper is only 66 $\mu$m, implying that current flows
predominantly near conductor surfaces at radio frequencies. This necessitates different
design strategies for RF circuits compared to DC circuits: hollow waveguides can
replace solid conductors, and surface plating (gold or silver) can reduce losses
without requiring solid precious metal construction.
\end{remark}

\begin{example}[Waveguide Mode Selection]
The WR-90 waveguide has a TE$_{10}$ cutoff at 6.56 GHz and TE$_{20}$ cutoff at
13.1 GHz. Operating in the X-band range (8.2--12.4 GHz) ensures single-mode propagation,
preventing modal dispersion and signal distortion. The operating frequency must satisfy
$1.25 f_{c,10} < f < 0.95 f_{c,20}$ to avoid cutoff and higher-order modes.
\end{example}

\subsection{Poynting Vector and Power Flow}

\begin{pycode}
# Calculate time-averaged Poynting vector magnitude
E_field_magnitude = 1.0  # V/m (normalized)
H_field_magnitude = E_field_magnitude / eta_0
power_density = 0.5 * E_field_magnitude * H_field_magnitude  # W/m^2

print(f"For a plane wave with $E_0 = 1$ V/m in free space, the time-averaged power density is $\\langle S \\rangle = E_0^2/(2\\eta_0) = {power_density*1e3:.3f}$ mW/m$^2$.")
\end{pycode}

The Poynting vector $\mathbf{S} = \mathbf{E} \times \mathbf{H}$ describes electromagnetic
energy flux. For plane waves, the time-averaged power density scales as $E_0^2$ and
inversely with medium impedance.

\section{Conclusions}

This comprehensive analysis of electromagnetic wave propagation demonstrates:

\begin{enumerate}
\item Fresnel coefficients predict reflectance and transmittance at dielectric interfaces,
with the air-glass Brewster angle at \py{f"{np.degrees(theta_b_air_glass):.1f}"}$^\circ$
and glass-air critical angle at \py{f"{np.degrees(theta_c_glass_air):.1f}"}$^\circ$
enabling total internal reflection.

\item Skin depth in copper decreases from 660 $\mu$m at 10 kHz to 2.1 $\mu$m at 10 GHz,
confining high-frequency currents to conductor surfaces and dictating RF circuit design.

\item Rectangular waveguides exhibit dispersive propagation above cutoff frequency, with
phase velocity exceeding $c$ while group velocity remains subluminal, satisfying
$v_p v_g = c^2$.

\item The dominant TE$_{10}$ mode in WR-90 waveguide has cutoff at
\py{f"{fc_10/1e9:.2f}"} GHz, enabling X-band operation (8.2--12.4 GHz) in single-mode regime.
\end{enumerate}

\section*{Further Reading}

\begin{thebibliography}{99}

\bibitem{balanis2012}
Balanis, C.A. \textit{Advanced Engineering Electromagnetics}, 2nd ed. Wiley, 2012.

\bibitem{jackson1999}
Jackson, J.D. \textit{Classical Electrodynamics}, 3rd ed. Wiley, 1999.

\bibitem{pozar2011}
Pozar, D.M. \textit{Microwave Engineering}, 4th ed. Wiley, 2011.

\bibitem{griffiths2017}
Griffiths, D.J. \textit{Introduction to Electrodynamics}, 4th ed. Cambridge University Press, 2017.

\bibitem{collin2001}
Collin, R.E. \textit{Foundations for Microwave Engineering}, 2nd ed. Wiley-IEEE Press, 2001.

\bibitem{harrington2001}
Harrington, R.F. \textit{Time-Harmonic Electromagnetic Fields}, 2nd ed. Wiley-IEEE Press, 2001.

\bibitem{born1999}
Born, M. and Wolf, E. \textit{Principles of Optics}, 7th ed. Cambridge University Press, 1999.

\bibitem{stratton2007}
Stratton, J.A. \textit{Electromagnetic Theory}. Wiley-IEEE Press, 2007.

\bibitem{elliott1993}
Elliott, R.S. \textit{Electromagnetics: History, Theory, and Applications}. IEEE Press, 1993.

\bibitem{kong2008}
Kong, J.A. \textit{Electromagnetic Wave Theory}. EMW Publishing, 2008.

\bibitem{cheng1989}
Cheng, D.K. \textit{Field and Wave Electromagnetics}, 2nd ed. Addison-Wesley, 1989.

\bibitem{ramo1994}
Ramo, S., Whinnery, J.R., and Van Duzer, T. \textit{Fields and Waves in Communication Electronics}, 3rd ed. Wiley, 1994.

\bibitem{marcuvitz1951}
Marcuvitz, N. \textit{Waveguide Handbook}. McGraw-Hill, 1951.

\bibitem{chen1983}
Chen, H.C. \textit{Theory of Electromagnetic Waves}. McGraw-Hill, 1983.

\bibitem{stutzman2012}
Stutzman, W.L. and Thiele, G.A. \textit{Antenna Theory and Design}, 3rd ed. Wiley, 2012.

\bibitem{orfanidis2016}
Orfanidis, S.J. \textit{Electromagnetic Waves and Antennas}. Rutgers University, 2016.

\bibitem{ulaby2015}
Ulaby, F.T. and Ravaioli, U. \textit{Fundamentals of Applied Electromagnetics}, 7th ed. Pearson, 2015.

\bibitem{johnk1988}
Johnk, C.T.A. \textit{Engineering Electromagnetic Fields and Waves}, 2nd ed. Wiley, 1988.

\end{thebibliography}

\end{document}
